% from ega2-1-fr.tex
\oldpage[II]{35}
à $(M\otimes_S N)_{(f)}$, on vérifie aussitôt que $\lambda_f$ et
$\lambda'_f$ sont des $S_{(f)}$-homomorphismes réciproques, d'où la
première partie de la proposition.

Pour démontrer la seconde partie, supposons que $M$ soit le conoyau d'un
homomorphisme $P\to Q$, de $S$-modules gradués, $P$ et $Q$ étant sommes
directes d'un nombre fini de modules de la forme $S(n)$~; utilisant
l'exactitude à gauche de $\operatorname{Hom}_S(L,N)$ en $L$ et l'exactitude
de $M_{(f)}$ en $M$, on est aussitôt ramené à prouver que $\mu_f$ est un
isomorphisme lorsque $M=S(n)$. Or, pour tout $z$ homogène dans $N$, soit
$u_z$ l'homomorphisme de $S(n)$ dans $N$ tel que $u_z(1)=z$~; on voit
aussitôt que $\eta:z\mapsto u_z$ est un isomorphisme de degré $0$ de
$N(-n)$ sur $\operatorname{Hom}_S(S(n),N)$. Il lui correspond un
isomorphisme
\[
  \eta_f:\bigl(N(-n)\bigr)_{(f)}
  \longrightarrow
  \bigl(\operatorname{Hom}_S(S(n),N)\bigr)_{(f)}.
\]
D'autre part, soit $\eta'_f$ l'isomorphisme
\[
  N_{(f)}\longrightarrow
  \operatorname{Hom}_{S_{(f)}}\bigl(S(n)_{(f)},N_{(f)}\bigr)
\]
qui, à tout $z'\in N_{(f)}$, fait correspondre l'homomorphisme $v_{z'}$ tel
que $v_{z'}(s/f^k)=sz'/f^{n+k}$ (pour
$s\in S_{n+k}=(S(n))_k$). On constate aisément que l'application composée
\[
  \bigl(N(-n)\bigr)_{(f)}
  \xrightarrow{\eta_f}
  \bigl(\operatorname{Hom}_S(S(n),N)\bigr)_{(f)}
  \xrightarrow{\mu_f}
  \operatorname{Hom}_{S_{(f)}}\bigl(S(n)_{(f)},N_{(f)}\bigr)
  \xrightarrow{{\eta'_f}^{-1}}N_{(f)}
\]
est l'isomorphisme $z/f^h\mapsto z/f^{h-n}$ de
$\bigl(N(-n)\bigr)_{(f)}$ sur $N_{(f)}$, donc $\mu_f$ est un isomorphisme.
\end{proof}

Si l'idéal $S_+$ est engendré par $S_1$, on déduit de
(\hyperref[II.2.5.13-fr]{2.5.13}) que pour tout idéal gradué $\mathfrak{J}$
de $S$ et tout $S$-module gradué $M$, on a
\[
\label{II.2.5.13.1-fr}
  \widetilde{\mathfrak{J}}\cdot\widetilde{M}
  =\widetilde{\mathfrak{J}\cdot M}
\tag{2.5.13.1}
\]
à un isomorphisme canonique près~; cela résulte en effet de la commutativité
du diagramme
\[
  \xymatrix{
    \widetilde{\mathfrak{J}}\otimes_{\mathcal{O}_X}\widetilde{M}
    \ar[rr]^\lambda\ar[dr]
    && \widetilde{\mathfrak{J}\otimes_S M}\ar[dl]\\
    & \widetilde{M} &
  }
\]
que l'on vérifie comme celle de
(\hyperref[II.2.5.11.3-fr]{2.5.11.3}).

\begin{corollary}[2.5.14]
\label{II.2.5.14-fr}
Supposons $S$ engendrée par $S_1$. Quels que soient $m$, $n$ dans
$\mathbf{Z}$, on a alors~:
\[
\label{II.2.5.14.1-fr}
  \mathcal{O}_X(m)\otimes_{\mathcal{O}_X}\mathcal{O}_X(n)
  =\mathcal{O}_X(m+n)
\tag{2.5.14.1}
\]
\[
\label{II.2.5.14.2-fr}
  \mathcal{O}_X(n)=\bigl(\mathcal{O}_X(1)\bigr)^{\otimes n}
\tag{2.5.14.2}
\]
à des isomorphismes canoniques près.
\end{corollary}

\begin{proof}
La première formule résulte de (\hyperref[II.2.5.13-fr]{2.5.13}) et de
l'existence de l'isomorphisme canonique de degré $0$,
$S(m)\otimes_S S(n)\xrightarrow{\sim}S(m+n)$, qui, à l'élément
$1\otimes 1$ (où le premier facteur $1$ est dans $(S(m))_{-m}$ et le second
dans $(S(n))_{-n}$) fait correspondre l'élément
$1\in(S(m+n))_{-m-n}$. Il suffit ensuite de démontrer la seconde formule
pour $n=-1$, et en vertu de (\hyperref[II.2.5.13-fr]{2.5.13}), cela revient
à voir que $\operatorname{Hom}_S(S(1),S)$ est canoniquement isomorphe à
$S(-1)$, vérification immédiate en remontant aux définitions
(\hyperref[II.2.1.2-fr]{2.1.2}) et en se souvenant que $S(1)$ est un
$S$-module monogène.
\end{proof}

